\documentclass[12pt,letterpaper]{article}
\usepackage{amsmath}
\usepackage{amsthm}
\usepackage{amsfonts}
\usepackage{amssymb}
\usepackage{amscd}
\usepackage{enumerate}
\usepackage{fancyhdr}
\usepackage{mathrsfs}
\usepackage{bbm}
\usepackage{framed}
\usepackage{mdframed}
\usepackage{cancel}
\usepackage{mathtools}
\usepackage{verbatim}
\usepackage{hyperref}
\usepackage[letterpaper,voffset=-.5in,bmargin=3cm,footskip=1cm]{geometry}
\setlength{\parindent}{0.0in}
\setlength{\parskip}{0.1in}
\allowdisplaybreaks
\headheight 15pt
\headsep 10pt
\newcommand\N{\mathbb N}
\newcommand\Z{\mathbb Z}
\newcommand\R{\mathbb R}
\newcommand\Q{\mathbb Q}
\newcommand\lcm{\operatorname{lcm}}
\newcommand\setbuilder[2]{\ensuremath{\left\{#1\;\middle|\;#2\right\}}}
\newcommand\E{\operatorname{E}}
\newcommand\V{\operatorname{V}}
\newcommand\Pow{\ensuremath{\operatorname{\mathcal{P}}}}

\DeclarePairedDelimiter\ceil{\lceil}{\rceil}
\DeclarePairedDelimiter\floor{\lfloor}{\rfloor}
\newcommand\hint[1]{\textbf{Hint}: #1}
\newcommand\note[1]{\textbf{Note}: #1}
\newcounter{enum22i}
\newenvironment{22enumerate}{\begin{enumerate}[a.]\itemsep0em\setcounter{enumi}{\value{enum22i}}}{\setcounter{enum22i}{\value{enumi}}\end{enumerate}}
\newenvironment{22itemize}{\begin{itemize}\itemsep0em}{\end{itemize}}
\fancypagestyle{firstpagestyle} {
  \renewcommand{\headrulewidth}{0pt}
  \lhead{\textbf{CSCI 0220}}
  \chead{\textbf{Discrete Structures and Probability}}
  \rhead{M. Littman}
}

\fancypagestyle{fancyplain} {
  \renewcommand{\headrulewidth}{0pt}
  \lhead{\textbf{CSCI 0220}}
  \chead{Homework 9}
  \rhead{\textit{ }}
}
\pagestyle{fancyplain}


\begin{document}
  \thispagestyle{firstpagestyle}
  \begin{center}
    {\huge \textbf{Homework 9}}

    {\large Due: \textit{Monday, July 26th at 11:59pm EDT}}
  \end{center}


     All homeworks are due the Monday after their release at 11:59pm EDT on Gradescope.

    Please do not include any identifying information about yourself in the handin, including your Banner ID.

    Be sure to fully explain your reasoning and show all work for full credit. We recommend checking out the sample proofs on the course website to see the level of rigor for which we're looking.
    
    

\subsection*{Problem 1}
{
{\setcounter{enum22i}{0}

For each of the following pairs of events, identify whether they are independent. Justify why or why not by using the fact that event $A$ and $B$ are independent if $P(A)=P(A|B)$.
\begin{22enumerate}
\item When flipping a fair coin three times:
\begin{22itemize}
  \item The first coin comes up tails ($T$)
  \item There is a run of exactly two heads
  (that is, there are two, but not three, heads flipped in a row)
\end{22itemize}



\item When generating a 0/1 string of length 5:
\begin{22itemize}
  \item The 3rd bit is a 1
  \item There are at least two 0s
\end{22itemize}
\end{22enumerate}
}
}

\subsection*{Problem 2}
{


{\setcounter{enum22i}{0}
 The 14 moons of Neptune work at Neptune's Intergalactic Inc. Each of them works at their respective desk. Triton, one of Neptune's worker moons, wants to prank Neptune, the boss of the company. Knowing Neptune hates disorder, Triton convinces everyone in the office to switch desks when Neptune is in its office. The prank is only successful if there is exactly one moon per desk after the jumble and the new arrangement is not identical to the previous one.
    \begin{22enumerate}
      \item If each moon chooses a desk uniformly and independently (potentially choosing their own desk), provide an expression for the probability the prank is successful (that is, each desk is assigned to exactly one moon and at least one moon has moved desks).
     
      \item Now, consider a case in which all desks are arranged in a circle, and the moons only have time to move to a desk directly to the left or right of their current desk. What is the probability the prank is successful? (Every moon moves desks this time.)
      
      \item Neptune is taking a little longer than usual to come out of its office---the moons will move to any desk exactly 2 to the right or left of their original position. In this case, what is the probability that the prank is successful?
    \end{22enumerate}
    
}
}


\subsection*{Problem 3}
{
{\setcounter{enum22i}{0}

Neptune is given the following equation:\[x_1 + x_2 + x_3 + x_4 =75\] where each $x_i$ must be non-negative.
	\begin{22enumerate}
	\item Count the number of solutions to this equation.\\
	\hint{Use Donuts and Separators!}
	\item Now suppose we require a solution with $x_1$ and $x_3$ strictly positive. Count the number of solutions under this new constraint.
	\item Let us now assume that each $x_i$ for $i\in[1,4]$ takes some value in the range $[0,75]$ uniformly at random. Given that $x_1$ and $x_3$ are strictly positive, what is the probability that $x_1 + x_2 + x_3 + x_4 = 75$?
	\end{22enumerate}


}
}

\subsection*{Mind Bender (Extra Credit)}
{

{\setcounter{enum22i}{0}

      Neptune is working on a research paper but its only local planetary library has made it so that each member can only take out a single book. In front of Neptune is an aisle of $n$ different books relevant to its paper, where $n$ is a positive integer.

      Neptune doesn't have enough time to go through every book, so it wants to select a book \textbf{uniformly at random}. In other words, Neptune looks to find a scheme that will guarantee that it selects the $i^{\text{th}}$ book with probability $\frac{1}{n}$, where $1 \leq i \leq n$.
      
      Neptune has some constraints. The planet is too tired to make multiple trips up and down the aisle, so it looks to go down the aisle only once. Furthermore, Neptune can only hold one book at a time.
      
      Come up with and prove a scheme that allows Neptune to walk away from the end of the aisle with the $i^{\text{th}}$ book with probability $\frac{1}{n}$.
      
}
}

  
\end{document}